\documentclass[11pt,a4paper]{article}

% --- Partitional primes: definition, provability, finiteness heuristic ---
% Compile with:  pdflatex PARTITIONAL.tex   (twice, for cross-references)

\usepackage[T1]{fontenc}
\usepackage[utf8]{inputenc}
\usepackage{lmodern}
\usepackage{amsmath,amssymb}
\usepackage{booktabs}
\usepackage{array}
\usepackage[table]{xcolor}
\usepackage{microtype}
\usepackage[margin=2.4cm]{geometry}
\usepackage{enumitem}
\usepackage[colorlinks=true,linkcolor=black,urlcolor=blue]{hyperref}

% The partitional operator: a left turnstile, like the one Oliver Lau chose (⊢).
\newcommand{\partop}{\vdash}
\newcommand{\Part}{\Pi}                 % the partitorial product
\newcommand{\jacobi}[2]{\left(\frac{#1}{#2}\right)}

\title{\textbf{Partitional Primes}\\[0.3em]
       \large Definition, provability, and a finiteness heuristic}
\author{Oliver Lau}
\date{}

\begin{document}
\maketitle

\noindent
A number class invented by Oliver Lau, together with the heuristic argument that
it almost certainly contains only \textbf{finitely many} primes. This note
records the definition, why these primes are cheap to prove, and the
probabilistic budget that predicts their scarcity --- so it does not get lost.

\section{Definition}

With $p(k)$ the integer \textbf{partition function}
\[
  p(0)=1,\; p(1)=1,\; p(2)=2,\; p(3)=3,\; p(4)=5,\; p(5)=7,\; p(6)=11,\;\dots
\]
define the \emph{partitorial}
\[
  \Part(n) \;=\; \prod_{k=0}^{n} p(k) \;=\; p(0)\cdot p(1)\cdot p(2)\cdots p(n),
\]
and the \textbf{Partitional numbers} (operator $\partop$, ASCII fallback \texttt{\S}):
\[
  n\partop \;=\; \Part(n) \pm 1 .
\]

\paragraph{Example ($n=10$).}
\[
  \Part(10)=1\cdot1\cdot2\cdot3\cdot5\cdot7\cdot11\cdot15\cdot22\cdot30\cdot42
          = 960\,498\,000,
\]
so the two candidates are $960\,498\,001$ and $960\,497\,999$.

\subsection*{How fast does \texorpdfstring{$\Part(n)$}{Pi(n)} grow?}

$\Part(n)$ grows \emph{much} faster than the factorial. The cleanest way to see
this is not to look at the number itself --- it quickly has thousands of digits
--- but at its \textbf{logarithm}. A handy fact: the number of decimal digits of
a whole number $X$ is about $\ln X / \ln 10 \approx \ln X / 2.303$. So ``how big
is $\Part(n)$'' really means ``how big is $\ln\Part(n)$,'' and that is the
quantity we track.

\paragraph{Step 1 --- turn the product into a sum.}
The logarithm turns multiplication into addition, which is exactly what we want
for a product:
\[
  \ln\Part(n) \;=\; \ln\!\prod_{k=0}^{n} p(k)
             \;=\; \sum_{k=0}^{n} \ln p(k).
\]
So the size of $\Part(n)$ is just the running total of the sizes of the
individual partition numbers $p(k)$.

\paragraph{Step 2 --- how big is a single partition number?}
Here we borrow one famous result, the \emph{Hardy--Ramanujan formula} (1918): the
partition numbers themselves already grow almost exponentially, and their
logarithm behaves like
\[
  \ln p(k) \;\sim\; \pi\sqrt{\tfrac{2}{3}\,k} \;\approx\; 2.565\,\sqrt{k}.
\]
In words: each partition number is roughly ``$e$ to the power of (a constant
times $\sqrt{k}$).'' This is the only outside ingredient we need --- everything
else is bookkeeping.

\paragraph{Step 3 --- add them all up.}
Summing Step 2 over $k=0,\dots,n$ means adding up $\sqrt{0}+\sqrt{1}+\dots+\sqrt{n}$.
A sum of $\sqrt{k}$ is well approximated by the area under the curve $\sqrt{x}$,
i.e.\ its integral:
\[
  \sum_{k\le n}\sqrt{k} \;\approx\; \int_{0}^{n}\!\sqrt{x}\,\mathrm{d}x
                       \;=\; \tfrac{2}{3}\,n^{3/2}.
\]
Putting Steps 2 and 3 together (the constant is
$\pi\sqrt{2/3}\cdot\tfrac23 \approx 2.565\cdot0.667 \approx 1.71$):
\[
  \boxed{\;\ln \Part(n)\;\sim\; 1.71\,\cdot n^{3/2}\;}
\]

\paragraph{What this means in plain terms.}
The exponent is $n^{3/2}=n\sqrt{n}$, i.e.\ ``$n$ times the square root of $n$.''
That one square-root factor is what makes Partitional numbers explode faster than
factorials: for the factorial, $\ln(n!)\sim n\ln n$ grows only like ``$n$ times
the (slow) logarithm of $n$,'' whereas here it grows like ``$n$ times the (much
larger) $\sqrt{n}$.'' Concretely, $\ln\Part(n)\approx 1.71\,n^{3/2}$ translates
via Step~1's digit rule ($\text{digits}\approx\ln\Part(n)/2.303$) into
\[
  \text{digits of }\Part(n)\;\approx\; 0.74\,\cdot n^{3/2}.
\]
Check it: at $n=316$ this predicts $0.74\cdot316^{3/2}\approx 4160$ --- the right
ballpark for the actual $3255$ digits (the formula is an asymptotic, a bit
generous at ``small'' $n$). At $n=1000$ it is already about $20\,000$ digits, and
at $n=10\,000$ roughly $740\,000$. This breakneck growth is precisely why
Partitional primes thin out so fast --- see the finiteness heuristic in
Section~\ref{sec:heuristic}.

\section{Why these primes are cheap to prove}

\paragraph{Two kinds of ``is it prime?''}
There is a fast, cheap test (a \emph{primality test}) that says ``this number is
almost certainly prime'' --- good enough to go hunting with, but not a proof. And
there is the harder job of a watertight \textbf{proof} that a number really is
prime, with no ``almost.'' For a number with thousands of digits, how expensive
that proof is depends entirely on \emph{what we know about the number's
neighbours}, $N-1$ and $N+1$.

\paragraph{The expensive route: ECPP.}
For a ``random'' large prime --- one with no special structure --- the only
general proof method is \textbf{ECPP} (Elliptic Curve Primality Proving). It
works, but it is slow and hungry for memory: proving a few-thousand-digit number
can take hours and gigabytes of RAM. This is the wall the rest of this project
keeps running into.

\paragraph{The cheap route: knowing the factors of a neighbour.}
There is a classical shortcut, more than a century old. \emph{If} you already
know the complete prime factorization of $N-1$ \emph{or} of $N+1$, then you can
\textbf{prove} $N$ prime quickly and with almost no memory. The catch is normally
the ``if'': factoring a big number is itself hopeless.

\paragraph{Why factoring $\Part(n)$ costs us nothing.}
The ``if'' is normally the whole problem --- but for Partitional numbers it is
free, and the reason is a small, almost childish observation: \textbf{the prime
factorization of a \emph{product} is just the factorizations of its pieces,
thrown together.} And $\Part(n)$ \emph{is} a product --- one we assembled
ourselves out of the small partition numbers $p(0),\dots,p(n)$. We never set a
factoring algorithm loose on the giant $\Part(n)$; we factor each little $p(k)$
(they have only a handful of digits) and collect. For $n=10$:
\[
  \Part(10)=2\cdot3\cdot5\cdot7\cdot11\cdot
            \underbrace{15}_{3\cdot5}\cdot
            \underbrace{22}_{2\cdot11}\cdot
            \underbrace{30}_{2\cdot3\cdot5}\cdot
            \underbrace{42}_{2\cdot3\cdot7}
          \;=\; 2^{4}\cdot3^{4}\cdot5^{3}\cdot7^{2}\cdot11^{2}
          \;=\; 960\,498\,000 .
\]
We have just read off the complete factorization of $960\,498\,000$ without ever
treating that nine-digit number as a monolith. The identical thing happens at
$n=316$: the product has $3255$ digits, yet its factorization is handed to us for
free, because every building block is tiny.

\paragraph{The trick, run in reverse.}
It is worth savouring how backwards this is. Normally you are handed a number $N$
and get stuck, because \emph{its} neighbours $N-1$ and $N+1$ both refuse to
factor. Here we go the other way round: we \emph{start} from a number we have
already factored completely --- $\Part(n)$, by construction --- and ask which
nearby numbers it now lets us certify. We are not picking a target and then
hunting for its factors; we are picking targets \emph{because} their factored
neighbour is already in hand. And a fully factored $\Part(n)$ has exactly two
immediate neighbours to offer: $\Part(n)+1$ and $\Part(n)-1$.

\paragraph{Two candidates, two directions.}
Because $\Part(n)$ sits right between those two candidates, the one free
factorization serves both of them --- once on each side. Take the first variant,
$N=\Part(n)+1$: its \emph{lower} neighbour is $N-1=\Part(n)$, precisely the
number we hold in factored form, and knowing the factors of $N-1$ is exactly what
the classical \textbf{$N-1$ test} wants (Pocklington's theorem, descended from
Lucas and later sharpened into the Brillhart--Lehmer--Selfridge, or ``BLS'',
form). The second variant, $N=\Part(n)-1$, tilts the same fact the other way: now
the factored number is the \emph{upper} neighbour, $N+1=\Part(n)$, which is the
input for the mirror-image \textbf{$N+1$ test} (built on Lucas sequences, likewise
in a BLS form). So the single pre-factored number $\Part(n)$ plays the part of
``$N-1$'' for one candidate and ``$N+1$'' for the other: two sides of the same
coin, and between them they cover both Partitional variants.

\subsection*{The idea behind the two tests, without formulae}

Both tests rest on one picture: a prime $N$ forces a certain \emph{cycle} to have
an exact length, and you could only have built a cycle of that length if $N$
really is prime. Checking the length precisely is where the factors come in.

\paragraph{The $N-1$ test.}
Picture a clock face whose positions are the numbers $1,2,\dots,N-1$. Pick a
starting value $a$ and ``tick'' by repeatedly multiplying by $a$, always wrapping
around (that is, keeping the remainder after dividing by $N$). Sooner or later
the hand lands back on its start; the number of ticks that takes is its
\emph{period}. A classical fact (Fermat's little theorem) says that when $N$ is
prime, some choice of $a$ runs with the \emph{full} period $N-1$ --- the hand
tours every position before returning home. Turn that around and it becomes a
proof: if you can \emph{exhibit} an $a$ whose period is exactly $N-1$, then the
clock truly has $N-1$ distinct positions, and that happens only when $N$ is
prime. The delicate part is confirming the period is the \emph{whole} $N-1$ and
not some shorter fraction of it --- and that is where you need the prime factors
of $N-1$: for each prime $q$ dividing $N-1$, you check that the hand is \emph{not}
yet home after only $(N-1)/q$ ticks. No factors of $N-1$, no check. A
pre-factored $N-1$ is the golden ticket.

\paragraph{The $N+1$ test.}
When it is $N+1$ you can factor instead, you play the very same game on a
\emph{larger} clock. To build it you widen the number system a little: roughly,
you admit one new ``imaginary'' quantity --- a square root that simply does not
exist among the ordinary remainders --- much as the complex numbers are made by
bolting $\sqrt{-1}$ onto the reals. In that enlarged world a prime $N$ conjures up
a clock of size $N+1$, and the role of ``keep multiplying by $a$'' is taken over
by a Fibonacci-like recipe (technically a \emph{Lucas sequence}: each term is a
fixed blend of the two before it). The rest is identical: exhibit such a sequence
whose period is the full $N+1$, confirm it does not collapse early at $(N+1)/q$
for each prime $q$ dividing $N+1$, and $N$ stands proven. Same idea, mirror side
--- and this time it is the factors of $N+1$ that you must know.

\paragraph{Why this is so cheap.}
Each check above is one ``run the hand to a certain number of ticks'' --- in
arithmetic terms, one modular exponentiation --- per prime factor of $\Part(n)$.
That is fast, and it uses only $O(N)$ memory (storage proportional to the
number's own length, nothing worse). It scales far beyond what ECPP can reach: we
proved $\Part(316)-1$ ($3255$ digits) deterministically in ${\sim}9$~minutes on a
laptop via the $N+1$ method, where plain ECPP on the same number would take hours
and heavy RAM.

\paragraph{The same trick in the wild.}
Partitional numbers are not the only ``product $\pm1$'' family, and the others are
hunted in earnest. The factorial primes $n!\pm1$, the primorial primes $p\#\pm1$,
and --- since 2025 --- the compositorial primes $n!/n\#\pm1$ (here $n\#$ is the
primorial, so $n!/n\#$ is the product of all \emph{composite} numbers up to $n$)
are all searched by the volunteer project \textbf{PrimeGrid}. Every record they
report is a genuine, fully \emph{proven} prime of hundreds of thousands of digits
--- not a mere probable prime --- and the reason is exactly the mechanism above:
each of these forms carries its own pre-factored neighbour, so the $N\pm1$ proof
stays cheap at sizes where ECPP would be hopeless. Partitional numbers share that
machinery precisely; what sets them apart is only their faster growth
(Section~\ref{sec:heuristic}), which alone pushes them onto the rare ``finitely
many'' side.

\medskip
\noindent\emph{(Fine print: for very large $n$ a few of the $p(k)$ themselves
become big and might not factor quickly. The BLS method still works as long as
the part we \emph{did} factor covers more than the cube root of $N$ --- which in
practice holds comfortably --- with ECPP needed only on whatever small leftover
remains. This has not been necessary in the range explored so far.)}

\section{Why there are probably only finitely many}\label{sec:heuristic}

This is not a proof (none is possible --- see Section~\ref{sec:proof}). It is a
back-of-the-envelope argument that nonetheless pins the answer down surprisingly
well.

\paragraph{Step 1 --- how likely is one candidate to be prime?}
A deep but standard result (the \emph{prime number theorem}) says: among numbers
near a large value $X$, roughly one in every $\ln X$ is prime. So if you hand me
a number of size $X$ and I know nothing special about it, my best guess is that
it is prime with probability about $1/\ln X$. For our candidate $\Part(n)\pm1$,
that size is $X=\Part(n)$, so
\[
  \Pr[\,\Part(n)\pm1 \text{ prime}\,] \;\approx\; \frac{C}{\ln\Part(n)} .
\]
The fudge factor $C$ just bundles two real effects: we get two shots at each $n$
(the $+1$ and the $-1$), and these candidates are known to have \emph{no small
factors} (they were built that way), which makes them a bit likelier to be prime
than a truly random number. We will pin $C$ down from the data in a moment.

\paragraph{Step 2 --- add up the chances over all $n$.}
The expected \emph{total} number of Partitional primes is just the sum of those
per-$n$ chances. And now the growth rate from Section~1 does all the work.
Because $\ln\Part(n)\approx 1.71\,n^{3/2}$, each term shrinks like $1/n^{3/2}$:
\[
  E_{\text{total}}
  \;=\; \sum_{n\ge2}\Pr[\,\Part(n)\pm1\text{ prime}\,]
  \;\approx\; C\cdot\sum_{n\ge2}\frac{2}{\ln\Part(n)}
  \;\propto\; \sum_{n\ge2}\frac{1}{n^{3/2}}.
\]

\paragraph{Step 3 --- does that sum add up to something finite?}
This is the crux, and it is a schoolbook distinction. An infinite sum of
ever-smaller terms can either pile up without limit (\emph{diverges}) or settle
on a finite total (\emph{converges}). The textbook dividing line is the exponent
on $n$: $\sum 1/n^{1}$ diverges (grows forever), while $\sum 1/n^{s}$ for any
exponent $s>1$ converges. Our exponent is $s=\tfrac32 > 1$, so
\[
  \sum_{n\ge2}\frac{1}{n^{3/2}} \;\;\text{\textbf{converges}} \;\;\Longrightarrow\;\;
  E_{\text{total}} \text{ is \textbf{finite}.}
\]
In plain words: the candidates grow so explosively (that $\sqrt{n}$ in the
exponent again) that the chance of being prime drops away fast enough for the
grand total of expected primes to be a finite number --- not infinity.

\paragraph{Step 4 --- put a number on it.}
We can fix $C$ by matching the prediction to reality. Up to $n=316$ we have found
$19$ primes, and the corresponding sum works out to
$2\sum_{n\le316}1/\ln\Part(n)=6.89$. Setting $C\cdot 6.89 = 19$ gives
$C\approx2.76$. Carrying that same $C$ out to \emph{all} $n$ (the tail of the sum
adds very little) gives the \textbf{total lifetime budget}:
\[
  E_{\text{total}}
  \;\approx\; 2.76\cdot 2\!\sum_{n\ge2}\frac{1}{\ln\Part(n)}
  \;\approx\; 2.76\cdot 7.05
  \;\approx\; \boxed{19.5\text{ primes --- over \emph{all} } n.}
\]

We have already found $19$. \textbf{The budget is essentially spent.} The table
below splits the remaining ${\sim}0.5$ of a prime across ranges of $n$ (``bands'').
$E$ is the expected number of hits in that band; $P(\ge1)=1-e^{-E}$ converts that
into the probability of finding \emph{at least one} (if hits land at random with
average count $E$, the chance of drawing a blank is $e^{-E}$, so the chance of
one or more is $1-e^{-E}$):

\begin{center}
\begin{tabular}{@{}lrr@{}}
\toprule
\textbf{Band} & \textbf{expected $E$} & \textbf{$P(\ge1)$}\\
\midrule
$(316,\;920]$      & $0.18$ & ${\sim}17\%$\\
$(920,\;1000]$     & $0.01$ & ${\sim}1\%$\\
$(1000,\;2000]$    & $0.07$ & ${\sim}7\%$\\
$(2000,\;4000]$    & $0.05$ & ${\sim}5\%$\\
$(4000,\;40000]$   & $0.07$ & ${\sim}7\%$\\
\midrule
\rowcolor{gray!15}
\textbf{beyond 920, over all $n$, forever} & $\mathbf{\approx 0.23}$ & $\mathbf{{\sim}21\%}$\\
\bottomrule
\end{tabular}
\end{center}

\noindent
So the expected number of Partitional primes \emph{beyond} the ones we already
have is well under one: \textbf{most likely there are no more, at any $n$ ---
ever.} Two honest caveats. First, $C$ was fitted from only $19$ data points at
small $n$, so read the individual percentages as rough orders of magnitude, not
precise odds. Second --- and this is the robust part --- the headline conclusion
(expected count much less than $1$ from here on) rides only on the growth rate
$\ln\Part(n)\sim n^{3/2}$, not on the exact value of $C$. Even if $C$ were off by
a factor of two, the answer would still be ``probably none left.''

\subsection*{Where this sits among ``generator'' primes}

It helps to compare Partitional numbers with the other classic recipes that build
candidates from a growing product or formula. The single thing that decides
``finitely many vs.\ infinitely many'' is whether the sum from Step~3 converges
--- which in turn is decided by how fast the numbers grow:

\begin{center}
\begin{tabular}{@{}lccc@{}}
\toprule
\textbf{Family} & \textbf{size $\ln(\text{value})$ grows like} & \textbf{the sum $\sum 1/\ln$} & \textbf{expectation}\\
\midrule
Primorial $p\#\pm1$          & $n\ln n$   & diverges & infinitely many\\
Factorial $n!\pm1$           & $n\ln n$   & diverges & infinitely many\\
Compositorial $n!/n\#\pm1$   & $n\ln n$   & diverges & infinitely many\\
\rowcolor{gray!15}
\textbf{Partitional $\Part(n)\pm1$} & $\mathbf{n^{3/2}}$ & \textbf{converges} & \textbf{finitely many}\\
Fermat $2^{2^k}+1$           & $2^k$      & converges (fast) & finitely many\\
\bottomrule
\end{tabular}
\end{center}

\noindent
The primorial ($p\#$, the product of all primes up to $p$), the factorial, and
the compositorial ($n!/n\#$, the product of all \emph{composite} numbers up to
$n$) all grow only moderately, their sums never settle, and all three are
expected to yield infinitely many primes. The Fermat numbers grow so savagely (a
tower of powers of $2$) that their sum converges almost instantly --- and indeed
only $5$ Fermat primes are known, with most experts expecting no more.
Partitional numbers land on the same ``finite'' side as Fermat, just with a more
generous budget (${\sim}19$ expected, versus the $5$ Fermat primes). That is the
company this new class keeps --- and, among the product-based families, it is the
lone outlier on the finite side.

\section{Can it be proven or disproven?}\label{sec:proof}

Short answer: \textbf{no --- not with today's mathematics}, in either direction.
The heuristic above tells us what to \emph{expect}; it does not settle anything.
This is the same frustrating situation as with the Fermat primes, where ``are
there finitely many?'' has been open for centuries.

\begin{itemize}[leftmargin=1.4em]
\item \textbf{We cannot prove there are finitely many.} To do that, you would
  need to show that from some point on, \emph{every} $\Part(n)\pm1$ is composite
  (non-prime). But there is no mechanism forcing that. The usual way a formula
  gets ``poisoned'' into always being composite is a fixed small divisor --- and
  these candidates are immune to exactly that. Here is why: the moment some prime
  $q$ shows up as a factor of one partition number $p(k)$, it is a factor of the
  whole product $\Part(n)$ for every $n\ge k$. So $\Part(n)$ is a multiple of
  $q$, which means $\Part(n)\pm1$ is \emph{one away} from a multiple of $q$ and
  therefore \emph{never} divisible by $q$. In other words the candidates are, by
  construction, coprime to (share no factor with) every prime that lives inside
  $\Part(n)$ --- they are designed to slip past all the small divisors. Ruling
  out primality forever would mean proving that an incredibly thin set of numbers
  eventually contains none at all, and no technique can do that.
\item \textbf{We cannot prove there are infinitely many either.} That is hard
  even for the ``easy'' families whose sums diverge (primorial, factorial), where
  everyone believes the primes go on forever but nobody can prove it. For a
  convergent family like ours it is only more hopeless.
\end{itemize}

\noindent
The one escape hatch that would turn this into a real theorem is an
\emph{algebraic} surprise: if $\Part(n)\pm1$ secretly always factored by a formula
--- the way $n^4+4$ always splits as $(n^2-2n+2)(n^2+2n+2)$ and so is composite
for every $n>1$ --- then finiteness (in fact, ``almost none'') would follow
instantly. Nothing hints at such a hidden factorization, and the coprimality
argument above makes a \emph{simple} one unlikely. But that identity, if it
existed, is the only known route to an actual proof.

\section{Known Partitional primes ($n\le320$, all proven)}

\begin{center}
\begin{tabular}{@{}rccc@{}}
\toprule
$n$ & variant & digits & method\\
\midrule
2   & $\Part(n)+1$ & 1    & $N-1$\\
3   & $\Part(n)+1$ & 1    & $N-1$\\
3   & $\Part(n)-1$ & 1    & $N+1$\\
4   & $\Part(n)+1$ & 2    & $N-1$\\
4   & $\Part(n)-1$ & 2    & $N+1$\\
5   & $\Part(n)+1$ & 3    & $N-1$\\
6   & $\Part(n)+1$ & 4    & $N-1$\\
6   & $\Part(n)-1$ & 4    & $N+1$\\
7   & $\Part(n)+1$ & 5    & $N-1$\\
7   & $\Part(n)-1$ & 5    & $N+1$\\
10  & $\Part(n)+1$ & 9    & $N-1$\\
13  & $\Part(n)+1$ & 15   & $N-1$\\
14  & $\Part(n)-1$ & 17   & $N+1$\\
18  & $\Part(n)+1$ & 27   & $N-1$\\
22  & $\Part(n)+1$ & 38   & $N-1$\\
67  & $\Part(n)+1$ & 258  & $N-1$\\
198 & $\Part(n)-1$ & 1535 & $N+1$\\
217 & $\Part(n)-1$ & 1780 & $N+1$\\
316 & $\Part(n)-1$ & 3255 & $N+1$\\
\bottomrule
\end{tabular}
\end{center}

\noindent
$19$ primes total ($11$ of form $\Part(n)+1$, $8$ of form $\Part(n)-1$). No
Partitional prime of either variant exists for $68\le n\le320$ except
$n=198,217,316$ --- consistent with the heuristic's thinning tail.

\section{Tools in this repo}

\begin{itemize}[leftmargin=1.4em]
\item \texttt{partitional.gp} --- explorer: \texttt{scan(nmax)} (table),
  \texttt{prove(n)} (both variants, deterministic, with timing),
  \texttt{report(nmax)} (list primes). The $N-1$ and $N+1$ BLS provers are
  validated against PARI's \texttt{isprime} for $n\le50$.
\item \texttt{hunt\_partitional.sh} --- resumable, \texttt{nice}-d PRP search
  over a range, appending hits to \texttt{results/partitional\_prp.log}; prove
  hits later with \texttt{gp -q partitional.gp} $\to$ \texttt{prove(n)}.
\end{itemize}

\end{document}
